\documentclass[11pt,a4paper]{scrartcl}

\usepackage[lithuanian]{babel}
\usepackage{../../1.\space Vienanariai\space ir\space daugianariai/manostilius_lt}
\usepackage{tasks}
\renewcommand{\theproblem}{\arabic{problem}}

\title{Vijeto teoremos taikymas}
\author{Andrius Berniukevičius}

\begin{document}

\maketitle

\section*{Kaip dirbti su uždaviniais?}

Pirmiausia kiekvieną uždavinį pabandykite išspręsti savarankiškai. Neskubėkite ieškoti užuominų — skirkite laiko apgalvoti, ką jau žinote ir kokį sprendimo būdą galėtumėte pritaikyti. Jei įstrigote, atsiverskite atitinkamą „Hintai 1“ užuominą. Jei jos nepakanka, tik tada skaitykite „Hintai 2“, o prireikus — „Hintai 3“. Užuominas naudokite palaipsniui, nežiūrėkite jų iš anksto.

Kiekvieno paties pagimtdytas sprendimas yra šimtą kartų vertingesnis už įvaikintą sprendimą.

\section{Uždaviniai}

\begin{problem}
Neišsprendę lygties $x^2 - 7x + 10 = 0$, raskite jos sprendinių sumą ir sandaugą.
\end{problem}


\begin{problem}
Neišsprendę lygties $x^2 + mx -2m-3 = 0$, raskite jos sprendinių sumą ir sandaugą.
\end{problem}

\begin{problem}
Sudarykite redukuotąją kvadratinę lygtį, kurios sprendiniai būtų $x_1 = 3$ ir $x_2 = -4$.
\end{problem}

\begin{problem}
Sudarykite redukuotąją kvadratinę lygtį, kurios sprendiniai būtų $x_1 = -5$ ir $x_2 = -1$.
\end{problem}

\begin{problem}
Sudarykite redukuotąją kvadratinę lygtį, kurios sprendiniai būtų $x_1 = \pi $ ir $x_2 = -2$.
\end{problem}

\begin{problem}
Sudarykite redukuotąją kvadratinę lygtį, kurios sprendiniai būtų $x_1 = 6- \sqrt{7} $ ir $x_2 = 6+ \sqrt{7} $.
\end{problem}

\begin{problem}
Sudarykite redukuotąją kvadratinę lygtį, kurios sprendiniai būtų $x_1 = -5-2 \sqrt{3} $ ir $x_2 = -5+2 \sqrt{3} $.
\end{problem}

\begin{problem}
Išspręskite lygčių sistemą: $\begin{cases} x + y = 10 \\ x - y = 4 \end{cases}$
\end{problem}

\begin{problem}
Išspręskite lygčių sistemą: $\begin{cases} 3x - y = 8 \\ x - y =  2 \end{cases}$
\end{problem}

\begin{problem}
Išspręskite lygčių sistemą: $\begin{cases} x \cdot y = 18 \\ x = 2y \end{cases}$
\end{problem}

\begin{problem}
Išspręskite lygčių sistemą: $\begin{cases} x \cdot y = 36 \\ y = 4x \end{cases}$
\end{problem}

\begin{problem}
Išspręskite lygčių sistemą: $\begin{cases} ab = 6 \\ a-b = 3{,}8 \end{cases}$
\end{problem}

\begin{problem}
Išspręskite lygčių sistemą: $\begin{cases} ab = 18 \\ a-b = 3 \end{cases}$
\end{problem}

\begin{problem}
Kvadratinės lygties $x^2 - 14x + c = 0$ sprendinių skirtumas yra lygus $4$. Susidarykite lygčių sistemą, raskite sprendinius ir apskaičiuokite $c$.
\end{problem}

\begin{problem}
Kvadratinės lygties $x^2 +bx +10 = 0$ sprendinių skirtumas yra lygus $1,5$. Susidarykite lygčių sistemą, raskite sprendinius ir koeficientą $b$.
\end{problem}

\begin{problem}
Lygties $x^2 + bx + 27 = 0$ vienas sprendinys yra tris kartus didesnis už kitą. Raskite teigiamą koeficiento $b$ reikšmę.
\end{problem}

\begin{problem}
Duota lygtis $x^2 - 10x + c = 0$. Žinoma, kad jos sprendinių kvadratų skirtumas $x_1^2 - x_2^2 = 40$. Raskite koeficientą $c$.
\end{problem}

\begin{problem}
Kvadratinės lygties $x^2 - 12x + c = 0$ sprendinių atvirkštinių dydžių suma yra lygi $\frac{3}{8}$ (t. y., $\frac{1}{x_1} + \frac{1}{x_2} = \frac{3}{8}$). Raskite lygties sprendinius ir koeficientą $c$.
\end{problem}

\begin{problem}
Kvadratinės lygties $x^2 + bx - 15 = 0$ vienas sprendinys yra $8$ vienetais didesnis už kitą. Raskite visas galimas koeficiento $b$ reikšmes.
\end{problem}

\begin{problem}
Lygties $x^2 - 11x + c = 0$ sprendiniai tenkina sąlygą $2x_1 - 3x_2 = 7$. Sudarykite lygčių sistemą, raskite lygties sprendinius ir $c$.
\end{problem}

\begin{problem}
Duota lygtis $2x^2 + bx + 18 = 0$. Jos sprendinių santykis yra $1 : 9$. Raskite koeficientą $b$, jeigu žinoma, kad $b < 0$.
\end{problem}

\begin{problem}
Duota lygtis $2x^2 + bx + 30 = 0$. Jos sprendinių santykis yra $3 : 5$. Raskite koeficientą $b$, jeigu žinoma, kad $b < 0$.
\end{problem}

\begin{problem}
Lygties $x^2 + bx + 20 = 0$ sprendinių kvadratų suma yra lygi $41$. Raskite abi galimas koeficiento $b$ reikšmes.
\end{problem}

\begin{problem}
Kvadratinės lygties $x^2 - 7x + c = 0$ sprendinių kubų suma yra lygi $133$. Raskite koeficientą $c$.
\end{problem}

\begin{problem}
Lygties $x^2 + bx + 10 = 0$ vienas sprendinys sudaro $40\%$ kito sprendinio. Raskite koeficientą $b$, jeigu žinoma, kad lygties sprendiniai yra teigiami skaičiai.
\end{problem}

\newpage
\section{Atsakymai}

\begin{enumerate}[label=\textbf{\arabic*.}, leftmargin=*]
    \item Sprendinių suma $7$, sandauga $10$.
    \item Sprendinių suma $-m$, sandauga $-2m-3$.
    \item $x^2+x-12=0$.
    \item $x^2+6x+5=0$.
    \item $x^2+(2-\pi)x-2\pi=0$.
    \item $x^2-12x+29=0$.
    \item $x^2+10x+13=0$.
    \item $(x,y)=(7,3)$.
    \item $(x,y)=(3,1)$.
    \item $(x,y)=(6,3)$ arba $(x,y)=(-6,-3)$.
    \item $(x,y)=(3,12)$ arba $(x,y)=(-3,-12)$.
    \item $(a,b)=\left(5,\frac{6}{5}\right)$ arba $(a,b)=\left(-\frac{6}{5},-5\right)$.
    \item $(a,b)=(6,3)$ arba $(a,b)=(-3,-6)$.
    \item Sprendiniai $5$ ir $9$, todėl $c=45$.
    \item $b=\frac{13}{2}$, sprendiniai $-4$ ir $-\frac52$; arba $b=-\frac{13}{2}$, sprendiniai $4$ ir $\frac52$.
    \item $b=12$.
    \item $c=21$.
    \item Sprendiniai $4$ ir $8$, todėl $c=32$.
    \item $b=-2$ arba $b=2$.
    \item Sprendiniai $x_1=8$, $x_2=3$, todėl $c=24$.
    \item $b=-20$.
    \item $b=-16$.
    \item $b=-9$ arba $b=9$.
    \item $c=10$.
    \item $b=-7$.
\end{enumerate}


\newpage
\section{Hintai 1}

\begin{enumerate}[label=\textbf{\arabic*.}, leftmargin=*]
    \item Pritaikykite Vijeto teoremą redukuotajai kvadratinei lygčiai.
    \item Pritaikykite Vijeto teoremą ir palikite atsakymą su parametru $m$.
    \item Žinodami lygties sprendinius, galite rasti redukuotosios lygties koeficientus.
    \item Žinodami lygties sprendinius, galite rasti redukuotosios lygties koeficientus.
    \item Žinodami lygties sprendinius, galite rasti redukuotosios lygties koeficientus.
    \item Žinodami lygties sprendinius, galite rasti redukuotosios lygties koeficientus.
    \item Žinodami lygties sprendinius, galite rasti redukuotosios lygties koeficientus.
    \item Sudėkite ir atimkite abi sistemos lygtis.
    \item Iš antrosios lygties išreikškite $x$ per $y$.
    \item Įstatykite $x=2y$ į sandaugos lygtį.
    \item Įstatykite $y=4x$ į sandaugos lygtį.
    \item Išreikškite $a$ per $b$ ir įstatykite į $ab=6$.
    \item Išreikškite $a$ per $b$ ir įstatykite į $ab=18$.
    \item Naudokite šaknų sumą ir jų skirtumą.
    \item Išreikškite šaknų skirtumo kvadratą per jų sumą ir sandaugą.
    \item Pažymėkite mažesnę šaknį $t$, tada kita yra $3t$.
    \item Išskaidykite kvadratų skirtumą dauginamaisiais.
    \item Paverskite atvirkštinių šaknų sumą viena trupmena.
    \item Išreikškite vieną šaknį per kitą, naudodami skirtumą $8$.
    \item Sudarykite sistemą iš šaknų sumos ir duoto tiesinio ryšio.
    \item Pažymėkite lygties sprendinius kaip $x_1=t$ ir $x_2= 9t$.
    \item Pažymėkite lygties sprendinius kaip $x_1=3t$ ir $x_2=5t$.
    \item Šaknų kvadratų sumai naudokite tapatybę $x_1^2+x_2^2=(x_1+x_2)^2-2x_1x_2$.
    \item Šaknų kubų sumai naudokite tapatybę $x_1^3+x_2^3=(x_1+x_2)^3-3x_1x_2(x_1+x_2)$.
    \item Pažymėkite didesnę teigiamą šaknį $t$, tada kita yra $0{,}4t$.
\end{enumerate}

\newpage
\section{Hintai 2}

\begin{enumerate}[label=\textbf{\arabic*.}, leftmargin=*]
    \item Jei šaknys yra $x_1,x_2$, tai $x_1+x_2=7$, o $x_1x_2=10$.
    \item $x_1+x_2=-m$, o $x_1x_2=-2m-3$.
    \item Sudaroma lygtis $x^2-(x_1+x_2)x+x_1x_2=0$.
    \item Sudaroma lygtis $x^2-(x_1+x_2)x+x_1x_2=0$.
    \item Naudokite $x^2-(x_1+x_2)x+x_1x_2=0$.
    \item Šaknų suma lygi $12$, o sandauga lygi $36-7$.
    \item Šaknų suma lygi $-10$, o sandauga lygi $25-(2\sqrt3)^2$.
    \item Sudėję lygtis gausite $2x=14$.
    \item Iš antrosios lygties gaunate $x=y+2$. Šią išraišką įstatykite į pirmąją lygtį.
    \item Gaunama $2y^2=18$.
    \item Gaunama $4x^2=36$.
    \item Gaunama kvadratinė lygtis $b^2+3{,}8b-6=0$.
    \item Gaunama kvadratinė lygtis $b^2+3b-18=0$.
    \item Šaknų suma lygi $14$, o jų skirtumas $4$.
    \item Šaknų skirtumo kvadratas lygus $(x_1+x_2)^2-4x_1x_2=b^2-40$.
    \item Pagal Vijeto teoremą $3t^2=27$.
    \item $x_1^2-x_2^2=(x_1-x_2)(x_1+x_2)$.
    \item $\frac1{x_1}+\frac1{x_2}=\frac{x_1+x_2}{x_1x_2}=\frac{12}{c}$.
    \item Pažymėkite šaknis $x_2=t$, $x_1=t+8$; jų sandauga lygi $-15$.
    \item Iš $x_1+x_2=11$ ir $2x_1-3x_2=7$ sudarykite dviejų lygčių sistemą.
    \item Šaknų sandauga lygi $\frac{18}{2}=9$, todėl $9t^2=9$.
    \item Pagal Vijeto teoremą šaknų sandauga lygi $15$, todėl $15t^2=15$. Šaknų sumą susiekite su koeficientu $b$.
    \item Šaknų sandauga lygi $20$, taigi $b^2-40=41$.
    \item Šaknų suma lygi $7$, todėl $x_1^3+x_2^3=7^3-21c$.
    \item Pažymėję didesnę šaknį $t$, o mažesnę $0{,}4t$, iš sandaugos gausite $0{,}4t^2=10$.
\end{enumerate}

\newpage
\section{Hintai 3}

\begin{enumerate}[label=\textbf{\arabic*.}, leftmargin=*]
    \item Atsakymas: suma $7$, sandauga $10$.
    \item Atsakymas: suma $-m$, sandauga $-2m-3$.
    \item Suma $-1$, sandauga $-12$; lygtis $x^2+x-12=0$.
    \item Suma $-6$, sandauga $5$; lygtis $x^2+6x+5=0$.
    \item Suma $\pi-2$, sandauga $-2\pi$; lygtis $x^2+(2-\pi)x-2\pi=0$.
    \item Suma $12$, sandauga $29$; lygtis $x^2-12x+29=0$.
    \item Suma $-10$, sandauga $13$; lygtis $x^2+10x+13=0$.
    \item $x=7$, $y=3$.
    \item Įstatę $x=y+2$ į pirmąją lygtį, gausite $3(y+2)-y=8$, taigi $y=1$ ir $x=3$.
    \item $y=\pm3$, o $x=2y$, taigi $(6,3)$ arba $(-6,-3)$.
    \item $x=\pm3$, o $y=4x$, taigi $(3,12)$ arba $(-3,-12)$.
    \item $(a,b)=\left(5,\frac65\right)$ arba $\left(-\frac65,-5\right)$.
    \item $(a,b)=(6,3)$ arba $(-3,-6)$.
    \item Šaknys $9$ ir $5$, todėl $c=9\cdot5=45$.
    \item $b^2-40=(1{,}5)^2$, todėl $b^2=\frac{169}{4}$ ir $b=\pm\frac{13}{2}$. Atitinkamos šaknys: $4$ ir $\frac52$, arba $-4$ ir $-\frac52$.
    \item $t=\pm3$; teigiama $b$ gaunama iš šaknų $-3$ ir $-9$: $b=12$.
    \item Šaknų skirtumas $4$, suma $10$, todėl šaknys $7$ ir $3$, o $c=21$.
    \item $12/c=3/8$, todėl $c=32$; lygties šaknys $4$ ir $8$.
    \item Šaknų suma gali būti $2$ arba $-2$, todėl $b=-2$ arba $b=2$.
    \item Iš sistemos gaunama $x_1=8$, $x_2=3$, todėl $c=24$.
    \item $t=\pm1$; sąlyga $b<0$ parenka šaknis $1$ ir $9$, todėl $b=-20$.
    \item Šaknys yra $3$ ir $5$, jų suma $8$, todėl $b=-2\cdot8=-16$.
    \item $b^2=81$, todėl $b=\pm9$.
    \item $343-21c=133$, todėl $c=10$.
    \item Teigiamos šaknys yra $5$ ir $2$, jų suma $7$, todėl $b=-7$.
\end{enumerate}

\end{document}
